\section{Allocation of a fixed total order}
\label{sec:allocation}

For positive real orders, define
\begin{equation}\label{eq:Fdef}
 F_{a,b}(z)=\sum_{n\ge1}\frac{H_{n-1}^{(b)}}{n^a}z^n,
 \qquad H_m^{(b)}=\sum_{j=1}^m j^{-b},\qquad
 g_{a,b}=\Im F_{a,b}(i).
\end{equation}
The series initially defines $F$ in the open unit disk. At $i$ we use its
principal analytic continuation. This convention matters: if $a+b=1$,
its boundary-series coefficients tend to $1/a$, and the ordinary series
does not converge. For integer orders our convention is
$F_{a,b}(z)=\Li_{a,b}(z,1)$, with decreasing summation indices.
Write $w=a+b$, and let $\beta$ and $\eta$ denote the Dirichlet beta and
eta functions. Their continuations are used whenever an ordinary series
does not apply.

The positive double integral in the incoming real-order report
\cite{realorder} is the starting point. We recall its short proof to
make the subsequent monotonicity argument independent of numerical evidence.

\begin{lemma}[Positive double integral, prior result]\label{lem:double}
For $a,b>0$ and $z\in\mathbb C\setminus[1,\infty)$,
\begin{equation}\label{eq:double}
 F_{a,b}(z)=\frac{z^2}{\Gamma(a)\Gamma(b)}
 \int_{0<y<x<1}
 \frac{(-\log x)^{a-1}\{\log(x/y)\}^{b-1}}
 {(1-xz)(1-yz)}\,dy\,dx.
\end{equation}
In particular $g_{a,b}<0$.
\end{lemma}
\begin{proof}
For nonnegative integers $j,k$, put $y=xv$ and integrate first in $v$
and then in $x$. The normalized moment of $x^jy^k$ is
$(j+k+2)^{-a}(k+1)^{-b}$. Expanding the denominators for $|z|<1$
therefore reproduces \eqref{eq:Fdef}. The weight has finite mass $2^{-a}$,
and both denominators stay uniformly away from zero on a compact subset
of the slit plane. This also proves holomorphic continuation there.
At $z=i$, minus the imaginary part of the rational factor is
$(x+y)/\{(1+x^2)(1+y^2)\}>0$.
\end{proof}

The following change of variables reveals a comparison which is not
stated in the inspected reports.

\begin{theorem}[Strict monotonicity at fixed total order]\label{thm:allocation}
For every $w>0$, the function
\[
 a\longmapsto -g_{a,w-a},\qquad 0<a<w,
\]
has strictly negative derivative. More precisely, let
$V_a\sim\operatorname{Beta}(a,w-a)$ and define
\begin{equation}\label{eq:Jw}
 J_w(v)=\frac1{\Gamma(w)}\int_0^\infty
 r^{w-1}\frac{e^{-vr}+e^{-r}}{4\cosh(vr)\cosh r}\,dr,
 \qquad 0\le v\le1.
\end{equation}
Then
\begin{align}
 -g_{a,w-a}&=\mathbb E J_w(V_a),\label{eq:betaallocation}\\
 \frac{d}{da}(-g_{a,w-a})
 &=\operatorname{Cov}\!\left(J_w(V_a),
          \log\frac{V_a}{1-V_a}\right)<0.\label{eq:covariance}
\end{align}
Its sharp endpoint bounds are
\begin{equation}\label{eq:sharpgaussian}
 \boxed{\quad
 \beta(w)-\beta(w-1)<-g_{a,w-a}
 <\frac{\beta(w)}2+\frac{\eta(w)}{2^{w+1}}.
 \quad}
\end{equation}
The lower endpoint is approached as $a\uparrow w$, and the upper endpoint
as $a\downarrow0$.
\end{theorem}
\begin{proof}
In \eqref{eq:double} put $x=e^{-s}$, $y=e^{-(s+t)}$. The positive
integrand for $-g$, including the Jacobian, becomes
\[
 \frac{s^{a-1}t^{b-1}}{\Gamma(a)\Gamma(b)}
 \frac{e^{-s}+e^{-(s+t)}}{4\cosh s\cosh(s+t)}\,ds\,dt.
\]
Set $r=s+t$ and $v=s/r$. The Jacobian $r$ combines with the powers
to give $r^{w-1}v^{a-1}(1-v)^{b-1}$. Dividing by the beta density
normalization proves \eqref{eq:betaallocation}. Every step is justified
by Tonelli's theorem for a positive integrand.

For fixed $r>0$, the function
\[
 v\longmapsto\frac{e^{-vr}+e^{-r}}{\cosh(vr)}
\]
is strictly decreasing. Its numerator decreases strictly and its
positive denominator increases. Consequently $J_w$ is strictly
decreasing on $[0,1]$. It is continuous there by domination by the
$v=0$ integrand, which is integrable at both endpoints of the $r$-axis.

Differentiate the beta density with $w$ fixed. Its score is
$\log(v/(1-v))-\psi(a)+\psi(w-a)$. Beta logarithmic moments are finite
for every interior parameter pair, so differentiation under the
bounded $J_w$ expectation is legitimate. Centering the score gives
\eqref{eq:covariance}. If $V'$ is an independent copy of $V$, then
\[
 2\operatorname{Cov}(f(V),h(V))
 =\mathbb E[(f(V)-f(V'))(h(V)-h(V'))].
\]
Here $f=J_w$ decreases strictly and $h(v)=\log(v/(1-v))$ increases
strictly. The product is strictly negative off the diagonal, which has
probability zero. This proves the derivative claim.

Since $\mathbb E V_a=a/w$, $V_a\to0$ in probability as $a\downarrow0$;
similarly $V_a\to1$ as $a\uparrow w$. Boundedness and continuity of
$J_w$ prove the asserted endpoint limits. At the first endpoint,
\begin{align}
 J_w(0)
 &=\frac1{\Gamma(w)}\int_0^\infty
 r^{w-1}\frac{1+e^{-r}}{4\cosh r}\,dr\notag\\
 &=\frac12\beta(w)+2^{-w-1}\eta(w).\label{eq:J0}
\end{align}
The beta and eta Mellin integrals converge for $w>0$; their series
identifications follow first where absolute convergence holds and then
through those integrals. For the other endpoint, put
$q(r)=(2\cosh r)^{-1}$. One has
\[
 \frac{e^{-r}}{2\cosh^2r}=q(r)+q'(r).
\]
Integration by parts for $w>1$ gives
$J_w(1)=\beta(w)-\beta(w-1)$. Both sides are holomorphic for
$\Re w>0$, so the same equality holds for all positive $w$.
For $0<w\le1$ it is an analytic-continuation identity, not a manipulation
of a convergent series for $\beta(w-1)$.
\end{proof}

This theorem concerns a redistribution of total order: increasing $a$
simultaneously decreases $b$. It is different from a comparison in one
index while the other remains fixed. The proof works below, on, and
above the finite signed-measure threshold $a+b=1$.

\section{The critical Euler conjecture is a theorem}
\label{sec:critical}

The incoming report \emph{The Sharp Real-Order Threshold for Harmonic
Polylogarithms} labels strict decrease of $-2g_{a,1-a}$ and its proposed
uniform Euler constant as Conjecture~\texttt{conj:constant}
\cite{realorder}. Theorem~\ref{thm:allocation} settles its monotonicity
part. We spell out the error implication because ordinary convergence
and the normalization of the remainder must be kept distinct.

For $a+b=1$ set
\begin{equation}\label{eq:eulerdef}
 f_m=\frac{H_{2m}^{(b)}}{(2m+1)^a},\qquad
 d_k=\sum_{j=0}^k(-1)^j\binom kj f_j,
 \qquad E_N=\sum_{k=0}^{N-1}\frac{d_k}{2^{k+1}}.
\end{equation}
The forward difference used here is $\Delta f_m=f_m-f_{m+1}$.

\begin{lemma}[Critical endpoint measure, prior result]\label{lem:criticalmeasure}
For $0<a<1$, $b=1-a$, there is a finite positive measure $\nu_{a,b}$
on $(0,1)$, with strictly positive density and total mass $1/a$, such that
\begin{equation}\label{eq:criticalmom}
 \frac{H_{n-1}^{(b)}}{n^a}
   =\frac1a-\int_0^1u^{n-1}\,d\nu_{a,b}(u).
\end{equation}
Equivalently, the signed representing measure is
$a^{-1}\delta_1-\nu_{a,b}$.
\end{lemma}
\begin{proof}
We recall the construction from \cite{realorder}. Write
$k_s(T)=T^{s-1}/\Gamma(s)$ and
\[
 h(t)=\frac1{1-e^{-t}}-\frac1t,\qquad
 C_{a,b}(T)=\int_0^T k_a(T-t)k_b(t)h(t)\,dt.
\]
The regularized Hurwitz integral \cite{dlmf-hurwitz} gives
\[
 \int_0^\infty e^{-nT}C_{a,b}(T)\,dT
 =n^{-a}\left(\zeta(b,n)-\frac{n^{1-b}}{b-1}\right).
\]
Since $b=1-a$, the last term contributes $1/a$. Define
\[
 d\nu_{a,b}(u)
 =\{-\zeta(b)k_a(-\log u)+C_{a,b}(-\log u)\}\,du.
\]
Both summands are positive: $\zeta(b)<0$ for $0<b<1$ and $h>0$.
The beta integral gives $0<C_{a,b}(T)<k_{a+b}(T)=1$.
Thus the density is integrable, and its $(n-1)$st moment is
$1/a-n^{-a}(\zeta(b)-\zeta(b,n))$. The finite Hurwitz difference is
$H_{n-1}^{(b)}$. Taking $n=1$ proves its mass and the formula.
\end{proof}

\begin{theorem}[Sharp critical-line Euler constant]\label{thm:criticalconstant}
For every $0<a<1$, $b=1-a$, and integer $N\ge1$,
\begin{equation}\label{eq:criticalerror}
 0<E_N-g_{a,b}<
 \left(\frac\pi4+\frac{\log2}{2}\right)2^{-N}.
\end{equation}
The displayed constant is optimal simultaneously over this entire
critical line and all $N\ge1$. For fixed $a$, the exact optimal
non-strict constant is $-2g_{a,1-a}$, attained at $N=1$, and
\begin{equation}\label{eq:criticalrange}
 \frac\pi2-1<-2g_{a,1-a}<\frac\pi4+\frac{\log2}{2}.
\end{equation}
The normalized errors $2^N(E_N-g_{a,b})$ decrease strictly to zero as
$N\to\infty$.
\end{theorem}
\begin{proof}
Equation~\eqref{eq:criticalmom} implies $d_0=0$ and, for $k\ge1$,
\[
 d_k=-\int_0^1(1-u^2)^k\,d\nu_{a,b}(u)<0.
\]
The finite measure justifies summing the absolutely convergent Euler
series under the integral. The analytic value at $i$ is its sum, and
the geometric tail is
\begin{equation}\label{eq:atomicerror}
 2^N(E_N-g_{a,b})=
 \int_0^1\frac{(1-u^2)^N}{1+u^2}\,d\nu_{a,b}(u).
\end{equation}
The integrand decreases strictly with $N$ for $0<u<1$ and converges
to zero. Its value at $N=1$ integrates to $-2g_{a,b}$ because $E_1=0$.
This proves the fixed-parameter statement and strict convergence.
Theorem~\ref{thm:allocation} at $w=1$ gives \eqref{eq:criticalrange},
using $\beta(1)=\pi/4$, $\beta(0)=1/2$, and $\eta(1)=\log2$.
It also gives the strict uniform upper bound in \eqref{eq:criticalerror}.
Finally, take $N=1$ and $a\downarrow0$ to show that no smaller uniform
constant can work.
\end{proof}

The atom at $u=1$ is essential in this argument. The ordinary terms
$(-1)^m f_m$ do not approach zero on the critical line; Euler acceleration
computes the analytic or Abel value. No divergent series has been
silently declared ordinarily convergent.

\section{A sharp global envelope and its unique maximum}
\label{sec:envelope}

Define
\begin{equation}\label{eq:Cdef}
 C(w)=\beta(w)+2^{-w}\eta(w),\qquad D(w)=C(w)-1\qquad(w>0).
\end{equation}
Theorem~\ref{thm:allocation} says that $C(w)$ is the sharp upper
envelope of $-2g_{a,b}$ at fixed $a+b=w$. The next theorem determines
its global shape. The normalized Mellin log-concavity principle used
in the proof is classical; see, for example, Proposition~3.2 of
Fradelizi--Li--Madiman \cite{flm}. We provide a strict version with
a direct proof rather than making a new priority claim for that principle.

\begin{lemma}[Strict normalized Mellin log-concavity]\label{lem:mellin}
Suppose $f:(0,\infty)\to(0,\infty)$ is strictly log-concave, all the
Mellin integrals and their first two parameter derivatives below are
finite, and these derivatives can be passed under the integral. Then
\[
 M(p)=\frac1{\Gamma(p)}\int_0^\infty r^{p-1}f(r)\,dr
\]
has $(\log M)''(p)<0$ for every $p>0$ in that domain.
\end{lemma}
\begin{proof}
Normalize $r^{p-1}f(r)$ to a probability density $P$. Choose a gamma
density $Q$ of shape $p$ and rate $\lambda$ with
\[
 \mathbb E_Q\log r=\psi(p)-\log\lambda=\mathbb E_P\log r.
\]
Both $\int(P-Q)$ and $\int\log r(P-Q)$ vanish. The logarithm of the
likelihood ratio is $\log f(r)+\lambda r+$ a constant, hence is
strictly concave. Its positive set is an interval. Normalization
forces both signs. A single crossing is impossible: multiplying the
difference by $\log(r/r_0)$ at that crossing would give a nonzero
quantity of one sign, contrary to the two matching moments. Thus
there are two crossings $r_1<r_2$ and the sign pattern of $P-Q$ is
negative, positive, negative.

Consequently
\[
 \int_0^\infty
 (\log r-\log r_1)(\log r-\log r_2)(P-Q)(r)\,dr<0.
\]
The constant and linear terms disappear by the matching moments.
Therefore $\operatorname{Var}_P(\log r)<
\operatorname{Var}_Q(\log r)=\psi_1(p)$. Differentiation of $\log M$
identifies its second derivative with precisely the difference of
these two variances. This proves the lemma.
\end{proof}

\begin{theorem}[Unique maximum of the sharp Gaussian envelope]
\label{thm:globalenvelope}
For every $w>0$, $D(w)>0$ and $(\log D)''(w)<0$. There is a unique
$w_*>0$ where $C$ attains its global maximum $C_*=C(w_*)$.
It satisfies $1<w_*<2$ and the following certified enclosures:
\begin{equation}
 1.30221658710124120959237170
 <w_*<1.30221658710124120959237171.\label{eq:wstar}
\end{equation}
Writing $L_C<C_*<U_C$, the certified value endpoints are
\begin{align}
 L_C&=1.136561103339509560952586375779942387681723540,\notag\\
 U_C&=1.136561103339509560952586375779942387681723541.
 \label{eq:Cstar}
\end{align}
In particular
\begin{equation}\label{eq:globalg}
 0<-2\Im F_{a,b}(i)<C_*\qquad(a,b>0),
\end{equation}
and this global constant is optimal, approached with $a\downarrow0$
and $b=w_*-a$.
\end{theorem}
\begin{proof}
Subtract the gamma integral for $1$ from \eqref{eq:J0}. One obtains
\begin{equation}\label{eq:Dmellin}
 D(w)=\frac1{\Gamma(w)}\int_0^\infty r^{w-1}f(r)\,dr,
 \qquad f(r)=\frac{e^{-r}(1-e^{-r})}{2\cosh r}>0.
\end{equation}
The logarithm of this density has
\begin{equation}\label{eq:logfcurvature}
 (\log f)''(r)=-\frac{e^r}{(e^r-1)^2}-\operatorname{sech}^2r<0.
\end{equation}
Near zero $f(r)=r/2+O(r^2)$, and at infinity it decays exponentially.
Thus all the conditions of Lemma~\ref{lem:mellin} hold on $w>0$.
This proves strict log-concavity and positivity.

The integral of $f(r)/r$ is finite, so $D(w)\to0$ as $w\downarrow0$.
Also $f(r)<e^{-2r}$, whence $D(w)<2^{-w}\to0$ at infinity. Positivity,
continuity, and these endpoint limits ensure an interior maximum.
Strict concavity of $\log D$ makes that maximum unique and nondegenerate.
The exact certificate in Section~\ref{sec:gausscertificate} proves
$C'(1)>0$, $C'(2)<0$, and the sharper derivative signs at the two
rational endpoints in \eqref{eq:wstar}. Strict log-concavity then
locates the unique critical point. The same certificate uses a tangent
upper bound for $\log D$ to enclose its value as in \eqref{eq:Cstar}.
Finally \eqref{eq:globalg} and optimality follow from the sharp allocation
endpoints in Theorem~\ref{thm:allocation}.
\end{proof}

\paragraph{Scope of the constant.}
The number $C_*$ is a bound for $-2g$, equivalently for the normalized
\emph{first} Euler remainder. This theorem does not assert that the same
constant bounds $2^N(E_N-g)$ for every $N$ throughout $a+b\ge1$.
The all-$N$ optimal constant proved here is the critical-line constant
in Theorem~\ref{thm:criticalconstant}. Off that line, normalized errors
need not decrease with $N$.

\begin{corollary}[Turan inequalities and a second normalization]
\label{cor:turan}
For $w>0$ and $0<h<w$,
\begin{equation}\label{eq:turan}
 D(w)^2>D(w-h)D(w+h).
\end{equation}
Moreover $2^wD(w)$ increases strictly from zero to one on $(0,\infty)$.
For example,
\begin{align}\label{eq:catalanineq}
 \left(G+\frac{\pi^2}{48}-1\right)^2
 >&\left(\frac\pi4+\frac{\log2}{2}-1\right)\notag\\
 &\hspace{4mm}\cdot\left(\frac{\pi^3}{32}
                         +\frac{3\zeta(3)}{32}-1\right),
\end{align}
where $G=\beta(2)$ is Catalan's constant.
\end{corollary}
\begin{proof}
The first assertion is strict midpoint log-concavity. For the second,
rewrite \eqref{eq:Dmellin} using a gamma random variable $R_w$ of shape
$w$ and rate $2$:
\[
 2^wD(w)=\mathbb E\frac{1-e^{-R_w}}{1+e^{-2R_w}}.
\]
The function of $R_w$ is strictly increasing from zero to one. If
$w_2>w_1$, couple $R_{w_2}=R_{w_1}+Y$ with an independent positive gamma
variable $Y$ of shape $w_2-w_1$ and rate $2$. This gives strict increase.
The limit at zero follows by the gamma mean, and the limit at infinity
by concentration or the gamma Laplace transform. The displayed
constant inequality is \eqref{eq:turan} at $w=2,h=1$.
\end{proof}

\section{All-order inversion of the envelope}
\label{sec:inversion}

There is a unique solution $W(x)>w_*$ of $C(W(x))=1+x$ for
$0<x<C_*-1$. Its large-order expansion is a concrete inverse
transseries for a combination of Dirichlet $L$-values. We include
an all-coefficient formula, since simply replacing $D(w)$ by $2^{-w}$
misses several noninteger scales.

For $n\ge3$, put
\begin{equation}\label{eq:lambdas}
 \epsilon_n=
 \begin{cases}1,&n\equiv1,2\pmod4,\\-1,&n\equiv0,3\pmod4,
 \end{cases}
 \qquad \lambda_n=\log_2(n/2)>0.
\end{equation}
For a finite-support multi-index $\mathbf k=(k_3,k_4,\ldots)$ of
nonnegative integers, write $K=\sum k_n$ and
$\Lambda(\mathbf k)=\sum k_n\lambda_n$. Set $(u)_0=1$ and
$(u)_j=u(u+1)\cdots(u+j-1)$.

\begin{theorem}[Explicit inverse generalized-power expansion]
\label{thm:inverse}
For every fixed $A>0$, as $x\downarrow0$,
\begin{equation}\label{eq:inverseall}
 W(x)=\frac{\log(1/x)}{\log2}
 +\frac1{\log2}
 \sum_{\substack{\mathbf k\ne0\\\Lambda(\mathbf k)<A}}
 \frac{(-1)^{K+1}(\Lambda(\mathbf k)+1)_{K-1}}
      {\prod_n k_n!}
 \left(\prod_n\epsilon_n^{k_n}\right)x^{\Lambda(\mathbf k)}
 +O(x^A).
\end{equation}
The sum is finite, and coefficients at coincident exponents are added.
Writing $\lambda=\log_2(3/2)$ and $\mu=\log_2(5/2)$, its first terms are
\begin{align}\label{eq:inversefirst}
 W(x)={}&\frac{\log(1/x)}{\log2}
 -\frac{x^\lambda+x}{\log2}
 -\frac{\lambda+1/2}{\log2}x^{2\lambda}
 +\frac{x^\mu}{\log2}\notag\\
 &-\frac{\lambda+1}{\log2}x^{\lambda+1}
 -\frac{(3\lambda+1)(3\lambda+2)}{6\log2}x^{3\lambda}
 +O\!\left(x^{\log_2(7/2)}\right).
\end{align}
\end{theorem}
\begin{proof}
For $w>1$ the absolutely convergent Dirichlet series gives
\[
 D(w)=2^{-w}+\sum_{n\ge3}\epsilon_n n^{-w}.
\]
Put $y=2^{-w}$ and $y=xt$. Then $D(w)=x$ is equivalent to
\begin{equation}\label{eq:timplicit}
 t+\sum_{n\ge3}\epsilon_n x^{\lambda_n}t^{1+\lambda_n}=1.
\end{equation}
For any fixed target order $A$, keep finitely many $n$ with
$\lambda_n\le A$ and regard the $x^{\lambda_n}$ as independent small
variables. The implicit derivative in $t$ at the origin is $1$.
The ordinary analytic implicit-function theorem therefore applies.

To justify discarding the rest, for a fixed sufficiently large $N$ and
$w\to\infty$ the integral test gives
\[
 \sum_{n\ge N}(2/n)^w
 \le(2/N)^w\left(1+\frac{N}{w-1}\right).
\]
Also $D(w)\sim2^{-w}$, so $y/x\to1$. Taking $\lambda_N>A$ makes
the discarded part $o(x^A)$, uniformly for $t$ in a fixed neighborhood
of $1$. Stability of the implicit solution preserves that error.
Since every $\lambda_n\ge\lambda_3>0$ and only finitely many are
below a fixed bound, the additive exponent semigroup is locally finite.
Taylor truncation hence gives an asymptotic expansion to every fixed order.

For the coefficients, introduce a scalar bookkeeping parameter $z$ and
write $t=1-z\sum_n\epsilon_n X_nt^{1+\lambda_n}$. Lagrange inversion
applied to $\log t$ gives, for $K\ge1$,
\[
 [X^{\mathbf k}]\log t
 =\frac{(-1)^K(K-1)!}{\prod k_n!}
   \left(\prod\epsilon_n^{k_n}\right)
   \binom{K+\Lambda-1}{K-1}.
\]
Now $W(x)=\log(1/x)/\log2-\log t/\log2$, and
$(K-1)!\binom{K+\Lambda-1}{K-1}=(\Lambda+1)_{K-1}$.
This proves \eqref{eq:inverseall}. In \eqref{eq:inversefirst}, the
exponent $\lambda+1=\lambda_6$ receives both the direct $n=6$ term
and the product of the $n=3$ and $n=4$ terms. Their combined coefficient
is $-(\lambda+1)/\log2$. This collision is why terms must be grouped
before an expansion is reported.
\end{proof}

\section{Exact enclosure of the Gaussian maximum}
\label{sec:gausscertificate}

The script \texttt{code/certify\_gaussian\_envelope.py} uses only
integer and rational arithmetic. A decimal grid with 90 digits is an
implementation choice for outward rounding, not an assumption about
the accuracy of a transcendental library. We give the error bounds
which connect its finite checks to Theorem~\ref{thm:globalenvelope}.

For $d=1,2$ let
\[
 A_d(w)=\sum_{j\ge0}\frac{(-1)^j}{(dj+1)^w},\qquad
 a_j(w)=(dj+1)^{-w}.
\]
Thus $A_1=\eta$, $A_2=\beta$. Their Euler partial sums are
\begin{equation}\label{eq:eulerAn}
 T_{d,N}(w)=2^{-N}\sum_{j=0}^{N-1}(-1)^j(dj+1)^{-w}
                \sum_{h=j+1}^N\binom Nh.
\end{equation}
With $T\sim\operatorname{Gamma}(w,1)$, the exact tail is
\begin{equation}\label{eq:Atail}
 A_d(w)-T_{d,N}(w)
 =2^{-N}\mathbb E\frac{(1-e^{-dT})^N}{1+e^{-dT}}.
\end{equation}
The gamma integral proves this identity first by a geometric expansion;
the bounded expectation then justifies it for every $w>0$.
It is positive and at most $2^{-N}$.

\begin{lemma}[Differentiated Euler enclosure]\label{lem:eulerderivative}
For $w\ge1$,
\[
 |A_d'(w)-T_{d,N}'(w)|<2^{1-N}.
\]
Consequently, if $C_N=T_{2,N}+2^{-w}T_{1,N}$, then
\begin{equation}\label{eq:Ccertbounds}
 0<C(w)-C_N(w)\le(1+2^{-w})2^{-N},\qquad
 |C'(w)-C_N'(w)|<4\,2^{-N}.
\end{equation}
\end{lemma}
\begin{proof}
Differentiating the gamma density in \eqref{eq:Atail} introduces
$\log T-\psi(w)$. Cauchy--Schwarz bounds its absolute mean by
$\sqrt{\psi_1(w)}\le\sqrt{\psi_1(1)}=\pi/\sqrt6<2$.
This proves the first bound. The product rule for $2^{-w}A_1$ gives
the derivative bound
$2^{-N}\{2(1+2^{-w})+2^{-w}\log2\}<4\,2^{-N}$.
All differentiated integrals are dominated on compact $w$-intervals.
\end{proof}

The finite centers in \eqref{eq:eulerAn} require logarithms and powers
at rational $w$. The verifier bounds $\log n$ by reducing to $1\le r\le2$
and using
\[
 \log r=2\sum_{j=0}^{M-1}\frac{v^{2j+1}}{2j+1}+R_M,
 \quad v=\frac{r-1}{r+1},\quad
 0\le R_M\le\frac{2v^{2M+1}}{(2M+1)(1-v^2)}.
\]
It bounds exponentials after scaling their argument to $[0,1]$, using
the positive Taylor series and a geometric tail, followed by outward
reciprocals and four squarings. Thus $(dj+1)^{-w}$ and
$-\log(dj+1)(dj+1)^{-w}$ both have rational enclosing intervals.
Every addition, multiplication and division rounds outward.

With $N=160$, the program proves positive and negative derivative
signs at the rational endpoints of \eqref{eq:wstar}. It separately proves
\[
 C'(1)>0.03260,\qquad C'(2)<-0.03561.
\]
To enclose the maximum value, let $\ell<u$ be the certified endpoints.
Strict log-concavity gives, for $\ell\le w\le u$,
\[
 D(w)\le D(\ell)
 \exp\!\left(\frac{C'(\ell)}{D(\ell)}(u-\ell)\right).
\]
Outward bounds for the right-hand side and the lower bound $C(\ell)$
give \eqref{eq:Cstar}. This is a rational certificate for the maximum,
not a floating-point optimizer followed by rounded digits.

Independent high-precision quadrature evaluates the beta-mixture and
Mellin integrals. In particular it checks the derivative at $w=1$
without differentiating the removable eta singularity numerically.
These diagnostics are useful checks on conventions, but the rational
certificate and the preceding analytic proofs supply the stated bounds.
